\begin{abstract}
Modern machine learning contains recovery thresholds, specialization,
trajectory changes, optimizer-dependent regimes, and collective alignment, but
these phenomena are usually studied with incompatible objects and evidential
standards. We present a numerical statistical-mechanics atlas that separates
seven layers---domain, state object, deformation, disorder, scale, observable,
and phenomenon---and records theory status, realism, and evidence on independent
axes. A machine-readable PhaseCard and one experiment grammar connect a
scientific question to an ensemble, finite-size path, intervention, estimator,
and falsifier.

Transformers form the deepest dossier: trainable synthetic experiments measure
attention, feed-forward, residual, objective, optimizer, scaling, and data
coordinates with independent outer-seed uncertainty and separate intensive
train, IID, and OOD risks. Frozen follow-ups test selected mechanisms without
turning exploratory screens into confirmatory evidence. Diffusion,
reinforcement learning, and multi-agent systems test whether the grammar remains
useful with trajectory, feedback, and population observables. Predictive phase
continuation is one atlas protocol: a calibrated surrogate must forecast an
untouched size or deformation against simple baselines, with censored boundaries
retained as outcomes. The contribution is a falsifiable map from questions to
artifacts and negative results. It does not assert a shared universality class,
nor that a smooth learning curve is a thermodynamic transition.
\end{abstract}
